%% LaTeX-Beamer template for KIT design
%% by Erik Burger, Christian Hammer
%% title picture by Klaus Krogmann
%%
%% version 2.0
%%
%% mostly compatible to KIT corporate design v2.0
%% http://intranet.kit.edu/gestaltungsrichtlinien.php
%%
%% Problems, bugs and comments to
%% burger@kit.edu

%\documentclass[18pt]{beamer}
\documentclass[18pt,handout, aspectratio=169]{beamer}
\usepackage{pgfpages}
%\setbeameroption{show notes on second screen = right}
%\documentclass[18pt,draft]{beamer}

\usetheme{kit2}
\setbeamercovered{invisible}
\setbeamertemplate{navigation symbols}{}


%%%% Zeichencodierung & Schrift
\usepackage[ansinew]{inputenc}
\usepackage[ngerman]{babel}
\usepackage{amssymb}
\let\raggedright\relax
\usepackage{ragged2e}
%\usepackage{color}


%%%% Einbinden von Graphiken
%\usepackage{graphicx}
\usepackage{graphics}
\usepackage{graphicx}
\usepackage{pgf,tikz}
%\pdfminorversion=6 

%%%% Bibliographie und Zitationen
\usepackage[round]{natbib}
\bibliographystyle{nachvorj}


%%%% Tabellen
\usepackage{multirow,booktabs,verbatim }
\setbeamertemplate{caption}[numbered]

%\graphicspath{{U:/Lehre/MethodenII/Grafiken/}}
%\graphicspath{{D:/uni/Lehre/MethodenII/Grafiken/}}
\graphicspath{{/Dropbox/Graphs/}}

\usepackage{hyperref}
\hypersetup{pdfpagemode=FullScreen}

\title{Datenauswertung}
\subtitle{Regression}
\author{Andreas Haupt}
\institute{Karlsruher Institut für Technologie}
\date{}
\def\sym#1{\ifmmode^{#1}\else\(^{#1}\)\fi}
\begin{document}

% change the following line to "ngerman" for German style date and logos
\selectlanguage{ngerman}

%title page
\begin{frame}
\titlepage
\end{frame}

\begin{frame}
\frametitle{Das Ziel: Isolation von Einflüssen für gute Vergleiche}
\small

Der Kern sozialwissenschaftlicher Fragen ist die Isolation von miteinander verbundenen Einflüssen. Es ist oft nicht sinnvoll, Gruppen ``naiv'' zu vergleichen, sondern nur diejenigen Elemente von Gruppen, die in anderen Eigenschaften identisch sind, um eine gemeinsame Vergleichsgrundlage zu haben. \vfill 

\begin{enumerate}
  \item Wie stark beeinflussen Frauen im Gegensatz zu Männern als Vorgesetze das Arbeitsklima? 
  \begin{enumerate}
    \item[] Wir wissen, dass es nicht zufällig ist, in welchen Organisationen Frauen Vorgesetzte sind. Die Art der Organisation beeinflusst daher sehr wahrscheinlich sowohl die Rekrutierung von Frauen in Führungspositionen als auch das Arbeitsklima. 
  \end{enumerate}
  \pause
  \item Führt die Einnahme von Vitamin D zu verringerten Gesundheitsproblemen? 
  \begin{enumerate}
    \item[] Wir wissen, dass Personen, die besonders gesundheitsbewusst sind, eher dazu neigen, Vitamin D zu nehmen und außerdem mehr auf ihre Gesundheit zu achten. 
  \end{enumerate}
  \pause
  \item Welche Bedeutung spielt Armut für den Verlust von Institutionenvertrauen? 
  \begin{enumerate}
    \item[] Personen in Armut sind überdurchschnittlich stärker von Arbeitslosigkeit betroffen. Arbeitslosigkeit selbst könnte Institutionenverlust hervorrufen oder ihn auch bestärken. 
  \end{enumerate}
\end{enumerate}

\end{frame}


\begin{frame}{Ziel}
\begin{columns}

\begin{column}{.5\linewidth}
\begin{center}
\resizebox{1\linewidth}{!}{\includegraphics{DAGS-LikeCats}}
\end{center}
\end{column}

\begin{column}{.5\linewidth}
Wir möchten den direkten Einfluss von (beliebig vielen) Eigenschaften/Treatments auf die Ausprägungen einer uns interessierenden Eigenschaft analysieren. Eigenschaften können aber auf viele Arten miteinander verbunden sein. Je nachdem, wie sie verbunden sind, ist es unterschiedlich sinnvoll, sie zu isolieren oder die Isolation muss auf unterschiedliche Art durchgeführt werden. 

\end{column}

\end{columns}
\end{frame}


\begin{frame}
\frametitle{Vier typische Konstellationen}
\begin{center}
\resizebox{1\linewidth}{!}{
\includegraphics{Leseleistung-Dritt1}
}
\end{center}
\end{frame}

\begin{frame}{Was ist ein Einfluss von X auf Y?}
Wir wollen wissen, ob ein höheres Leseinteresse zu einer besseren Leseleistung führt. 

\begin{center}
\resizebox{.7\linewidth}{!}{
\includegraphics{lehre-scatter-lesen-lesint-empty}
}
\end{center}

\end{frame}

\begin{frame}
\frametitle{Einflüsse $\neq$ Korrelation}
Bisher besprochene \textit{Zusammenhangsmaße} geben uns \textbf{nicht} an, wie stark der Einfluss einer Eigenschaft auf eine andere Eigenschaft ist. \vfill
\begin{center}
\includegraphics[width=0.80\textwidth]{Correlation_examples}
\end{center}
\end{frame}

\begin{frame}
\frametitle{Ziel - in statistical terms}

Wir möchten ein statistisches Verfahren, 

\begin{enumerate}
	\item das uns die erwartete Veränderung der Ausprägung einer Eigenschaft angibt, wenn sich in einer erklärenden Eigenschaft \textbf{eine} Ausprägung ändert; 
	\item das es erlaubt, so viele Drittvariableneinflüsse wie nötig zu kontrollieren;
  \item das es uns ermöglicht, abzuschätzen, wie sicher wir uns über unsere Ergebnisse sein können.
\end{enumerate}
\end{frame}


\begin{frame}
\frametitle{Die Erklärung von Unterschieden des Outcomes}
\begin{center}
\resizebox{.8\linewidth}{!}{
\includegraphics{Leseleistung-Dritt1}
}
\end{center}
\small
Im folgenden behaupten wir, dass sich die Leseleistung jede:r individuellen Schüler:in \textbf{additiv} aus mehreren Komponenten zusammensetzt: 

\begin{enumerate}
\item die Leseleistung einer Basisgruppe, von der aus wir vergleichen;
\item \textit{systematischen Einflüssen} auf die Leseleistung nach spezifischen Ausprägungen von Eigenschaften (= erklärende Merkmale/Versuchsgruppen);
\item \textit{zufälligen Unterschieden} in der Leseleistung.  
\end{enumerate}

\end{frame}


\frame{
\frametitle{Die Regressionsgleichung}
Diese Behauptung können wir mathematisch auf folgende Art und Weise ausdrücken: 
\[
\underbrace{y_i}_{\substack{\text{individuelle} \\ \text{Leseleistung}}} = \underbrace{\hat{\alpha}_{\color{white}{i}}}_{\substack{\text{Leseleistung} \\ \text{Basisgruppe}}} +  \underbrace{\hat{\beta}\cdot x_i}_{\substack{\text{systematische} \\ \text{Gruppenunterschiede}}}  + \underbrace{\epsilon_i}_{\substack{\text{individuelle} \\ \text{Abweichung}}}
\]
\vfill
\begin{tabular}{llp{.8\linewidth}} \toprule 
$y_i$ & & Gemessene Werte der Leseleistung pro Person $i$\\
$\hat{\alpha}$ & ``alpha'' & Erwarteter Wert der Leseleistung für die Basisgruppe\footnote{Dieser Wert wird typischerweise als "`Konstante"' bezeichnet.} \\
$\hat{\beta}x_i$ & ``beta'' & Veränderung im erwarteten Wert der Leseleistung, wenn sich die Ausprägung der Gruppe um eine Einheit ändert\\ 
$\epsilon_i$ & ``epsilon'' & (zufällige) Abweichung der Leseleistung der Person vom erwarteten Leseleistungswert der Gruppe, in der diese Person ist \\ \bottomrule 
\end{tabular}

}

\frame{
\frametitle{Regression = fancy Gruppenvergleich}

Unser Ziel ist es, gute Vergleiche auf Gruppenebene anzustellen. Individuelle Werte sind nicht von Belang. Daher ignorieren wir bis auf Weiteres \textit{individuelle Abweichungen}, die wir als Residuen ($\epsilon_i$) bezeichnet haben. Unsere Regressionsgleichung verändert sich dann zu: 

\[
\underbrace{\hat{y}_i}_{\substack{\text{geschätzte} \\ \text{Leseleistung}}} = \underbrace{\hat{\alpha}_{\color{white}{i}}}_{\substack{\text{Leseleistung} \\ \text{Basisgruppe}}} +  \underbrace{\hat{\beta}\cdot x_i}_{\substack{\text{systematische} \\ \text{Gruppenunterschiede}}}  
\]
\vfill
Wir vergleichen immer Gruppen von Personen miteinander, die in der Variablen $X$ unterschiedliche Ausprägungen haben. Diese Variable kann das Geschlecht, den Schultyp oder das Leseinteresse messen. Wenn wir Gruppen auf diese Weise vergleichen, können wir eine Aussage machen, wie groß die Unterschiede zwischen den Gruppen \textit{typischerweise} sind. Diese Unterschiede drückt $\hat{\beta}$ aus. 

}


\frame{
\frametitle{Schätzungen von Einflüssen}
\small

Eine Möglichkeit, einen Einfluss zu schätzen, wäre ein ``educated guess''. Wir könnten uns auf Erfahrungswerte stützen oder einfach raten und das dann mit den Daten abgleichen. Ein ``educated guess'' könnte sein: 

\begin{itemize}
  \item Schüler:innen ohne jegliches Leseinteresse erwarten wir eine Leseleistung von 50 Punkten.
  \item Wir erwarten, dass sich die Leseleistung um 2 Punkte für jeden Skalenpunkt des Leseinteresses erhöht. 
\end{itemize}

\begin{center}
\begin{tabular}{lccccccc} \toprule 
$\hat{\alpha}$ & &	$\hat{\beta}$	& & Leseinteresse & &	   $\hat{y}$\\ \midrule
50 &+&  2 & $\cdot$ & 0  & = & 50\\
50 &+&  2 & $\cdot$ & 1  & = & 52\\
50 &+&  2 & $\cdot$ & 2  & = & 54\\
$\cdots$ & \\ 
50 &+&	2 & $\cdot$ & 10  &	= &	70\\
50 &+&	2 & $\cdot$ & 100 &	= &	250\\
50 &+&	2 & $\cdot$ & 200 &	= &	450\\
 \toprule
\end{tabular}
\end{center}
}


\begin{frame}
\frametitle{Unterschiede in X führen zu Unterschieden in Y}

\begin{center}
\resizebox{.82\linewidth}{!}{
  \includegraphics{lehre-ss20-scatter-lesen-lesint-reg0}
}
\end{center}

\end{frame}




\frame{
\frametitle{Was ist eine gute Schätzung?}

\begin{center}
  \begin{minipage}[b]{.5\linewidth} % [b] => Ausrichtung an \caption

\resizebox{1.1\linewidth}{!}{
\includegraphics{lehre-ss20-scatter-lesen-lesint-reg1}
}

  \end{minipage}
  \hspace{.05\linewidth}% Abstand zwischen Bilder
  \begin{minipage}[b]{.4\linewidth} % [b] => Ausrichtung an \caption 
Unser ``educated guess'' erscheint jedoch problematisch. 
\begin{itemize}
	\item Unsere Erwartungen erscheinen systematisch zu niedrig. 
	\item Es liegen deutlich mehr positive als negative Abweichungen der gemessenen Werte von der Schätzung vor. 
	\item Wir unterschätzen daher den Einfluss des Leseinteresses auf die Leseleistung. 
\end{itemize}
\end{minipage}
\end{center}

}


\frame{
\frametitle{Abweichungen der beobachteten von den erwarteten Werten (=Residuen)}

\begin{center}
  \begin{minipage}[b]{.5\linewidth} % [b] => Ausrichtung an \caption
\resizebox{1.1\linewidth}{!}{
\includegraphics{lehre-ss20-scatter-lesen-lesint-reg1}
}
\quad \\[2em]  
  \end{minipage}
  \hspace{.05\linewidth}% Abstand zwischen Bilder
  \begin{minipage}[b]{.4\linewidth} % [b] => Ausrichtung an \caption 

\begin{align*}
  \varepsilon_i &= y_i - \hat{y}_i\\ &= y_i - (\hat{\alpha}+\hat{\beta}x_i)
\end{align*}

Vorhersagefehler (Residuen) sind der Abstand zwischen den \textit{tatsächlichen} und den \textit{geschätzten} Werten. 
\begin{align*} 
396.1 - (50 + 2 \cdot 55.1) &=  235.9\\
216.1 - (50 + 2 \cdot 173) &= -179.9
\end{align*}
\quad \\[2em]  

\end{minipage}
\end{center}

}


\frame{
\frametitle{Was ist eine gute Schätzung? - Residuen}

Wir behaupten: Diejenige Schätzung ist am Besten, die die kleinst mögliche Summe der (quadrierten) Abstände zwischen den tatsächlichen und den geschätzten Werten generiert. 


\begin{itemize}
	\item Die Residuen werden quadriert, weil sich negative und positive Abweichungen von der Schätzung nicht wegsummieren sollen. Wir möchten das ``Gesamtpaket '' aller Residuen haben. 
	\item Schätzungen, die hohe einzelne Abweichungen verursachen, erhalten durch die Quadrierung ebenfalls eine schlechtere Bewertung, weil hohe quadrierte Abweichungen die Gesamtsumme aller quadrierten Residuen stark beeinflussen. 
\end{itemize}

}



\frame{
\frametitle{Wie minimiert man die Summe der quadrierten Residuen?}

\begin{align*}
\mathrm{min} =& \sum\limits_{i=1}^n(y_i - \hat{y}_i)^2 \\ 
			=&\sum\limits_{i=1}^n(y_i - (\hat{\alpha} + \hat{\beta}x_i))^2 
\end{align*}

Mathematisch lösbar durch: Ableitung nach $\hat{\alpha}$ und $\hat{\beta}$ und diese beiden ersten Ableitungen auf 0 setzen und nach $\hat{\alpha}$ oder $\hat{\beta}$ auflösen. 

}


\frame{
\frametitle{\textbf{O}rdinary-\textbf{L}east-\textbf{S}quares-Koeffizienten}

Daraus erhält man (für den Fall von \textbf{einer} genutzten Variablen): 

\begin{equation*}
  \hat{\alpha}=\bar{y}-\hat{\beta} \bar{x} 
\end{equation*}

Die Schätzung von $\hat{\alpha}$ verdeutlicht, dass dies ein Mittelwert (=Erwartungswert!) von $y$ ist, aus dem alle anderen durch $x$ bedingten systematischen Unterschiede herausgerechnet wurden. 

\begin{equation*}
\hat{\beta} = \dfrac{\sum(x_i-\bar{x})(y_i - \bar{y})}{\sum(x_i-\bar{x})^2} = \dfrac{\sum x_iy_i - n\cdot\bar{x}\cdot\bar{y}}{\sum x_i^2 - n\cdot \bar{x}^2}
\end{equation*}

Der $\hat{\beta}$-Koeffizient drückt ein Verhältnis der Covarianz von $x$ und $y$ in Relation zur Variation der $x$-Variable aus. Oder einfacher: Welche Unterschiede im Zusammenhang von $x$ und $y$ beruhen auf der Variation von $x$? 

}


\frame{
\frametitle{\textbf{O}rdinary-\textbf{L}east-\textbf{S}quares-Koeffizienten}

Für den Fall des Leseinteresses brauchen wir folgende Angaben, um beide Koeffizienten zu berechnen: 

\begin{center}
\begin{tabular}{ccccccc} \toprule
$\bar{y}$ & $\bar{x}$ & $\sum x_iy_i$ & $\sum x_i^2$ & n\\ \midrule
 303.08 &  111.35 & 35192353 & 13845040 & 1000\\ \bottomrule
\end{tabular}
\end{center}

\[ \hat{\beta} = \dfrac{\sum x_iy_i - n\cdot\bar{x}\cdot\bar{y}}{\sum x_i^2 - n\cdot \bar{x}^2} = \dfrac{35192353 - 1000 \cdot 303.08 \cdot 111.35 }{ 13845040 - 1000 \cdot 111.35^2} \approx 1. \]
Für jede Erhöhung einer Einheit des Leseinteresses erhöht sich die Leseleistung durchschnittlich um 1 Einheit. 

\[ \hat{\alpha} = \bar{y}-\hat{\beta} \bar{x} = 303.08 - 1 \cdot 111.35 = 191.73 \]
Schüler:innen ohne Leseinteresse haben eine durchschnittliche Leseleistung von 191.73. 

}


\begin{frame}
\frametitle{Schätzung des Einflusses nach OLS}

\begin{center}
\resizebox{.82\linewidth}{!}{
  \includegraphics{lehre-ss20-scatter-lesen-lesint-reg3}
}
\end{center}

\end{frame}

\begin{frame}
\frametitle{Schätzung des Einflusses nach OLS}

\begin{center}
\resizebox{.82\linewidth}{!}{
  \includegraphics{lehre-ss20-scatter-lesen-lesint-reg4}
}
\end{center}

\end{frame}



\end{document}



\frame{
  \frametitle{Regression mit Dummy-Variablen}

Der Einfluss des Leseinteresses auf die Leseleistung bezieht sich auf zwei metrisch gemessene Merkmale. \\ \vfill Wie können wir aber einen solchen Einfluss schätzen, wenn wir eine nominale oder ordinale erklärende Variable nutzen wollen? 

}


\frame{
  \frametitle{Regression mit Dummy-Variablen}

  \begin{center}
  \resizebox{.7\linewidth}{!}{
    \includegraphics{lehre-ss20-scatter-lesen-proj}
  }
  \end{center}

}

\frame{
\frametitle{\textbf{O}rdinary-\textbf{L}east-\textbf{S}quares-Koeffizienten}

Für den Fall der Projektteilnahme brauchen wir folgende Angaben, um beide Koeffizienten zu berechnen: 

\begin{center}
\begin{tabular}{ccccccc} \toprule
$\bar{y}$ & $\bar{x}$ & $\sum x_iy_i$ & $\sum x_i^2$ & n\\ \midrule
 303.08 &  0.25 & 85533.233 & 250 & 1000\\ \bottomrule
\end{tabular}
\end{center}

\[ \hat{\beta} = \dfrac{\sum x_iy_i - n\cdot\bar{x}\cdot\bar{y}}{\sum x_i^2 - n\cdot \bar{x}^2} = \dfrac{85533.233  - 1000 \cdot 303.08 \cdot 0.25 }{ 250 - 1000 \cdot 0.25^2} \approx 52. \]
Wenn sich die Projektteilnahme um eine Einheit erhöht (0 $\rightarrow$ 1),  erhöht sich die Leseleistung durchschnittlich um 52 Einheiten. 

\[ \hat{\alpha} = \bar{y}-\hat{\beta} \bar{x} = 303.08 - 52 \cdot 0.25 = 290.1 \]
Schüler:innen, die nicht am Projekt teilgenommen haben, haben eine durchschnittliche Leseleistung von 290.1. 

}


\begin{frame}[fragile]
\frametitle{OLS-Regression - PC-Version}
\scriptsize
\begin{verbatim}
      Source |       SS           df       MS      Number of obs   =     1,000
-------------+----------------------------------   F(1, 998)       =    392.61
       Model |  1442165.73         1  1442165.73   Prob > F        =    0.0000
    Residual |   3665943.2       998  3673.28978   R-squared       =    0.2823
-------------+----------------------------------   Adj R-squared   =    0.2816
       Total |  5108108.93       999  5113.22215   Root MSE        =    60.608

------------------------------------------------------------------------------
       lesen |      Coef.   Std. Err.      t    P>|t|     [95% Conf. Interval]
-------------+----------------------------------------------------------------
      lesint |   .9986589   .0504008    19.81   0.000     .8997553    1.097563
       _cons |    191.879   5.930404    32.36   0.000     180.2415    203.5165
------------------------------------------------------------------------------

\end{verbatim}
\end{frame}


\begin{frame}
  \frametitle{Regression - Zeitschriftversion}

\begin{center}
{

\begin{tabular}{l*{1}{c}}
\toprule
                    &\multicolumn{1}{c}{(1)}\\
                    &\multicolumn{1}{c}{Leseleistung}\\
\midrule
Leseinteresse       &        1.00\sym{***}\\
                    &      (0.05)         \\
\addlinespace
Konstante           &      191.88\sym{***}\\
                    &      (5.93)         \\
\midrule
N                   &     1000         \\
$R^2$                 &        0.28         \\
\bottomrule
\multicolumn{2}{l}{\footnotesize Standardfehler in Klammern}\\
\multicolumn{2}{l}{\footnotesize \sym{*} \(p<0.05\), \sym{**} \(p<0.01\), \sym{***} \(p<0.001\)}\\
\end{tabular}
}
\end{center}
\end{frame}


\begin{frame}[fragile]
\frametitle{The big picture}
\begin{center}
\resizebox{1\linewidth}{!}{
{
\def\sym#1{\ifmmode^{#1}\else\(^{#1}\)\fi}
\begin{tabular}{l*{11}{c}}
\toprule  
                   & & \multicolumn{2}{c}{Additiver Effekt} & \multicolumn{2}{c}{Mediation} & \multicolumn{2}{c}{Konfundierung} & \multicolumn{3}{c}{Interaktion}\\
                   & Wahrer Einfluss &\multicolumn{1}{c}{(1)}&\multicolumn{1}{c}{(2)}&\multicolumn{1}{c}{(3)}&\multicolumn{1}{c}{(4)}&\multicolumn{1}{c}{(5)}&\multicolumn{1}{c}{(6)}&\multicolumn{1}{c}{(7)}&\multicolumn{1}{c}{(8)}&\multicolumn{1}{c}{(9)}\\
\midrule
Ältere(s) Geschwister & 25&       26.38\sym{***}&       26.52\sym{***}&       25.60\sym{***}&       26.23\sym{***}&       26.63\sym{***}&       26.44\sym{***}&       26.49\sym{***}&       27.24\sym{***}&       27.31\sym{***}\\
Humanistische Ausrichtung& 15 &                     &       15.64\sym{***}&       15.11\sym{***}&       15.17\sym{***}&       15.10\sym{***}&       13.92\sym{***}&       13.72\sym{***}&       15.71\sym{***}&       15.97\sym{***}\\ \midrule
Teilhabe am Projekt & 0 &                    &                     &       50.30\sym{***}&        1.45         &        1.26         &       -1.27         &       -1.28         &       -2.50         &       -2.20         \\
Leseinteresse       & 1 &                    &                     &                     &        0.97\sym{***}&        0.97\sym{***}&        0.99\sym{***}&        0.99\sym{***}&        1.00\sym{***}&        1.00\sym{***}\\ \midrule
Nachhilfe           & 40 &                    &                     &                     &                     &       23.97\sym{***}&       39.74\sym{***}&       40.12\sym{***}&       40.32\sym{***}&       40.54\sym{***}\\
Leseschwäche        & -80 &                    &                     &                     &                     &                     &      -79.86\sym{***}&      -79.29\sym{***}&      -77.34\sym{***}&      -77.95\sym{***}\\ \midrule
\multicolumn{3}{l}{Bildung der Eltern (Referenz: mittlere Bildung)}\\
Niedrig             & -30 &                    &                     &                     &                     &                     &                     &      -34.91\sym{***}&      -33.08\sym{***}&      -27.87\sym{***}\\
Hoch                &  40 &                   &                     &                     &                     &                     &                     &       24.08\sym{***}&       25.42\sym{***}&       41.01\sym{***}\\ \addlinespace
Frau                & 70 &                    &                     &                     &                     &                     &                     &                     &       87.63\sym{***}&       70.56\sym{***}\\ \addlinespace
Frau \& Eltern niedr. Bildung& 20 &                     &                     &                     &                     &                     &                     &                     &                     &       21.18\sym{***}\\
Frau \& Eltern hohe Bildung&  30 &                  &                     &                     &                     &                     &                     &                     &                     &       31.23\sym{***}\\ \midrule
Konstante            & (120) &      301.82\sym{***}&      292.34\sym{***}&      279.93\sym{***}&      183.49\sym{***}&      175.06\sym{***}&      176.64\sym{***}&      172.29\sym{***}&      125.11\sym{***}&      117.99\sym{***}\\
\midrule
$R^2$                &  &        0.02         &        0.03         &        0.11         &        0.24         &        0.27         &        0.35         &        0.41         &        0.73         &        0.73         \\
N                  & &     5000         &     5000         &     5000         &     5000         &     5000         &     5000         &     5000         &     5000         &     5000         \\
\bottomrule
\multicolumn{10}{l}{\footnotesize \sym{*} \(p<0.05\), \sym{**} \(p<0.01\), \sym{***} \(p<0.001\)}\\
\end{tabular}
}
}
\end{center}
\end{frame}



\frame{
\frametitle{Berechnung der Güte des Modells}

Forscher:innen können daran interessiert sein, welchen Anteil die Unterschiede der erklärenden Variablen an den Unterschieden im Outcome haben. D.\,h.: Wie stark können wir Unterschiede in Y auf Unterschiede in X zurückführen?

\begin{equation}
\begin{split}
  R^2 &= \frac{\sum\limits_{i=1}^n(y_i-\bar{y})^2-\sum\limits_{i=1}^n(y_i-\hat{y_i})^2}{\sum\limits_{i=1}^n(y_i-\bar{y})^2}\\
   &= \frac{\sum\limits_{i=1}^n(\hat{y_i}-\bar{y})^2}{\sum\limits_{i=1}^n(y_i-\bar{y})^2} = \frac{\mbox{SQE}}{\mbox{SQT}}
\end{split}
\end{equation}

}

\frame{
\frametitle{Berechnung der Güte des Modells}
Genauso wie bei einer ANOVA berechnen wir die Abweichungen zum Gesamtmittelwert und die Abweichungen zum Vorhersagewert für jede Einheit der erklärenden Variable. Für den Fall des Leseinteresses ergibt sich daher: 
  
\begin{center} 
\begin{tabular}{lll}
$\sum (y_i-\bar{y})^2$ =&  5108108.93     &   \\ 
$\sum (\hat{y}_i-\bar{y}_i)^2$ =& 1442165.73  &    \\ 
\end{tabular}
\end{center}

\[
  R^2 = \frac{1442165.73  }{5108108.93  } \approx 0.28
\]



Unterschiede im Leseinteresse \textit{erklären} ca. 28\% der Unterschiede der Leseleistung.\footnote{Das bedeutet \textbf{nicht}, dass wir 28\% aller Leseleistungen korrekt vorhersagen oder dass 28\% aller Fälle auf der Regressionsgerade liegen würden. } 
}




\frame{
\frametitle{Wie kann man die Signifikanz der Regressionsparameter bestimmen?}

Problem: Die OLS-Koefizienten werden aufgrund der Stichproben-Werte bestimmt. Man möchte allerdings eine Aussage über den \textit{Effekt} in der Grundgesamtheit machen. 

Daher muss durch Signifikanztests überprüft werden, ob die Regressionsparameter sich von 0 unterscheiden. Ist dies nicht der Fall, dann können wir nicht mehr ausschließen dass es \textit{keinen} Effekt (Effekt = 0) gibt. 

}

\frame{
\frametitle{Wie kann man die Signifikanz der Regressionsparameter bestimmen?}

Prinzipiell sind diese Signifikanztest spezielle t-Tests. Sie prüfen die $H_0$, dass die Regressionsparameter 0 sind. 

\[
t = \frac{\mbox{geschätzter Parameter} - 0}{\mbox{Standardfehler des Parameters}}
\]

\begin{equation}
  t_{\hat{\alpha}} = \frac{\hat{\alpha}}{\hat{\sigma}_{\hat{\alpha}}} \approx t_{df}
\end{equation}

\begin{equation}
  t_{\hat{\beta}} = \frac{\hat{\beta}}{\hat{\sigma}_{\hat{\beta}}} \approx t_{df}
\end{equation}

}

\frame{
\frametitle{Wie kann man die Signifikanz der Regressionsparameter bestimmen?}

Die Standardfehler der $\beta$-Koeffizienten für den Fall einer erklärenden Variable berechnen sich durch: 

\begin{equation}
\begin{split}
  \hat{\sigma}_{\hat{\beta}} &= \sqrt{\frac{\frac{\sum\limits_{i=1}^n(y_i-\hat{y}_i)^2}{n-2}}{\sum\limits_{i=1}^n(x_i-\bar{x})^2}}\\
                             &= \sqrt{\frac{\mbox{Varianz der Residuen}}{\mbox{Variation der unab. Variable}}}
\end{split}
\end{equation}

}

\frame{
\frametitle{Wie kann man die Signifikanz der Regressionsparameter bestimmen?}

Die Standardfehler der $\alpha$-Koeffizienten berechnen sich durch: 

\begin{equation}
 \hat{\sigma}_{\hat{\alpha}} = \sqrt{\frac{\sum\limits_{i=1}^n(y_i-\hat{y}_i)^2}{n-2} \cdot \left(  \frac{1}{n} + \frac{\bar{x}^2}{\sum\limits_{i=1}^n(x_i-\bar{x})^2}  \right)} 
\end{equation}

}



\frame{
\frametitle{Wie kann man die Signifikanz der Regressionsparameter bestimmen?}

\[\hat{\sigma}_{\hat{\beta}} = \sqrt{\frac{\frac{3665943.2}{1000-2}}{1446041.5} } \approx 0.05 \]

Wir erwarten eine typische Variation von $\hat{\beta}$  durch Zufallsschwankungen aus alternativen Stichproben von 0.05 Einheiten. 

\[ \hat{\sigma}_{\hat{\alpha}} = \sqrt{ \frac{3665943.2}{1000-2} \cdot \left(  \frac{1}{1000} + \frac{111.35^2 }{1446041.5}  \right)  } = 5.93 \] 

Wir erwarten eine typische Variation von $\hat{\alpha}$  durch Zufallsschwankungen aus alternativen Stichproben von 5.93 Einheiten. 
}

\frame{
\frametitle{Wie kann man die Signifikanz der Regressionsparameter bestimmen?}

\[ t_{\hat{\beta}} = \frac{1}{0.05} = 20 \]

Bei einer Irrtumswahrscheinlichkeit von 5\% lauten die kritischen Werte für eine t-Verteilung mit 998 Freiheitsgraden (ungerichtet / gerichtet ): 1.977 / 1.656. 

\vfill

Wir können daher die $H_0$ verwerfen, dass es in der Grundgesamtheit keinen Einfluss des Leseinteresses auf die Leseleistung gibt. 

}


\frame{
\frametitle{Wie kann man die Signifikanz der Regressionsparameter bestimmen?}

Mit den nun zur Verfügung stehenden Informationen können wir auch das Konfidenzintervall für $\hat{\beta}$ berechnen: 

\[ 
OG / UG = 1  \pm 1.977 \cdot 0.05 = 1.099 / 0.90
\]

Bei einer Vielzahl von alternativen, unabhängigen Zufallsstichproben würden wir bei 95\% aller Stichproben $\hat{\beta}$-Koeffizienten zwischen 0.90 und 1.099 erwarten. Wir gehen davon aus, dass dieses Intervall den Populationswert überdeckt. Diese Annahme ist mit einer Wahrscheinlichkeit von 5\% falsch.

}






\begin{frame}
\frametitle{\textbf{O}rdinary-\textbf{L}east-\textbf{S}quares-Koeffizienten}
\begin{center}
  \begin{minipage}[b]{.4\linewidth} % [b] => Ausrichtung an \caption
\resizebox{1.1\linewidth}{!}{
\includegraphics{ss19-lehre-reg-dis}
}
\vfill 
  \end{minipage}
  \hspace{.05\linewidth}% Abstand zwischen Bilder
  \begin{minipage}[b]{.5\linewidth} % [b] => Ausrichtung an \caption 
Der $\beta$-Koeffizient drückt Unterschiede der geschätzten Werte pro Ausprägung von $X$ aus. In unserem Fall haben wir nur zwei Ausprägungen $X = 0$ und $X = 1$. Es sind aber unendlich viele Ausprägungen denkbar. Die Logik bleibt aber immer die gleiche: Wenn $X$ um eine Einheit steigt, ändert sich $y$ um $\beta$-Einheiten. 
\vfill 
Der $\alpha$-Koeffizient drückt immer unsere Schätzung aus, wenn alle im Modell befindlichen Einflüsse nicht vorhanden sind, d.h., wenn alle X-Werte 0 sind. 
\end{minipage}
\end{center}
\end{frame}

\begin{frame}
\frametitle{\textbf{O}rdinary-\textbf{L}east-\textbf{S}quares-Koeffizienten}

Bisher haben wir die Korrelation der erklärenden Variablen untereinander nicht einbezogen. Wir wissen aber, dass einige unserer erklärenden Eigenschaften untereinander stark korreliert sind. Zum Beispiel wissen wir, dass es keinen direkten Einfluss der Elternschaft auf die Spendenhöhe gibt, sondern dass dies eine Mediation über das Spielplatzprojekttreatment ist (Folie 2). Wenn wir, wie bisher, einen $\hat{\beta}$-Koeffizient berechnen, kommen wir aber zu einem anderen Schluss: 


\begin{center}
\begin{tabular}{ccccccc} \toprule
$\bar{y}$ & $\bar{x}$ & $\sum x_iy_i$ & $\sum x_1^2$ & n\\ \midrule
42.787 & 0.686 & 15270.627 & 343 & 500\\ \bottomrule
\end{tabular}
\end{center}

\[ \hat{\beta} = \dfrac{\sum x_iy_i - n\cdot\bar{x}\cdot\bar{y}}{\sum x_i^2 - n\cdot \bar{x}^2} = \dfrac{15270.627 - 500 \cdot 42.787 \cdot 0.686 }{ 343  - 500 \cdot 0.686^2} = 5.52. \]

\end{frame}





\frame{
\frametitle{Multivariate OLS}

Im Fall von zwei erklärenden Variablen ändern sich die Berechnungen der Koeffizienten folgendermaßen:\footnote{Die Logik für mehr als zwei erklärende Variablen ist gleich, allerdings erhöht sich der Grad der notwendigen mathematischen Komplexität, auf den ich hier verzichte.} 

\begin{equation}
  \hat{\alpha}= \bar{y}-\hat{\beta}_1\bar{x}_1-\hat{\beta}_2\bar{x}_2
\end{equation}

\begin{equation}
  \hat{\beta}_1=\frac{r_{yx_1}-r_{yx_2}\cdot r_{x_1x_2}}{1-r_{x_1x_2}^2}\cdot \frac{s_y}{s_{x_1}}
\end{equation}

\begin{equation}
  \hat{\beta}_2=\frac{r_{yx_2}-r_{yx_1}\cdot r_{x_1x_2}}{1-r_{x_1x_2}^2}\cdot \frac{s_y}{s_{x_2}}
\end{equation}


}

\frame{
\frametitle{Multivariate OLS}

Die Schätzung des Einflusses von Elternschaft ($X_1$) auf die Spendenhöhe ($Y$) unter Kontrolle des Einflusses von Spielplatzprojekttreatments ($X_2$) benötigt daher folgende Informationen: 

\begin{center}
\begin{tabular}{ccccccc} \toprule
$r_{yx{_1}}$ & $r_{yx_2}$ & $r_{x_1x_2}$ & $s_y$ & $s_{x_1}$ & $s_{x_2}$ \\ \midrule
0.2414 & 0.3916  & 0.4506 & 10.631 & 0.465 & 0.498 \\ \bottomrule
\end{tabular}
\end{center}

\begin{equation*}
  \hat{\beta}_1=\frac{r_{yx_1}-r_{yx_2}\cdot r_{x_1x_2}}{1-r_{x_1x_2}^2}\cdot \frac{s_y}{s_{x_1}} = \frac{0.2414 - 0.3916\cdot 0.4506}{1-0.4506^2}\cdot \frac{10.631}{0.465} \approx 1.86.
\end{equation*}

\begin{equation*}
  \hat{\beta}_2=\frac{r_{yx_2}-r_{yx_1}\cdot r_{x_1x_2}}{1-r_{x_1x_2}^2}\cdot \frac{s_y}{s_{x_2}} = \frac{0.3916  - 0.2414\cdot 0.4506}{1-0.4506^2}\cdot \frac{10.631}{0.498 } \approx 7.59.
\end{equation*}

}


\begin{frame}[fragile]
\frametitle{Multivariate OLS-Regression - PC-Version}
\scriptsize
\begin{verbatim}

      Source |       SS           df       MS      Number of obs   =       500
-------------+----------------------------------   F(2, 497)       =     46.85
       Model |  8946.11409         2  4473.05705   Prob > F        =    0.0000
    Residual |  47451.2513       497  95.4753547   R-squared       =    0.1586
-------------+----------------------------------   Adj R-squared   =    0.1552
       Total |  56397.3654       499  113.020772   Root MSE        =    9.7711

------------------------------------------------------------------------------
     spenden |      Coef.   Std. Err.      t    P>|t|     [95% Conf. Interval]
-------------+----------------------------------------------------------------
      eltern |   1.863997   1.054641     1.77   0.078    -.2081069      3.9361
     spielpl |   7.573047   .9834943     7.70   0.000     5.640727    9.505366
       _cons |   37.35786    .808385    46.21   0.000     35.76958    38.94613
------------------------------------------------------------------------------
\end{verbatim}
\end{frame}


\frame{
  \frametitle{Kompakte Publikationsdarstellung}

\begin{center}
\resizebox{.7\linewidth}{!}{
  \input{lehre-ss19-multreg.tex}
}
\end{center}

}




\frame{
\frametitle{Berechnung der Güte des Modells}
\begin{equation}
\begin{split}
  R^2 &= \frac{\sum\limits_{i=1}^n(y_i-\bar{y})^2-\sum\limits_{i=1}^n(y_i-\hat{y_i})^2}{\sum\limits_{i=1}^n(y_i-\bar{y})^2}\\
   &= \frac{\sum\limits_{i=1}^n(\hat{y_i}-\bar{y})^2}{\sum\limits_{i=1}^n(y_i-\bar{y})^2} 
\end{split}
\end{equation}

}
